\originalchapter{热辐射与光子气体}\label{ch:radiation}
\originalsection{平衡辐射与辐射通量}
同温腔体若通过滤光片相连, 任何频率上出现净能流都会破坏平衡. 所以大腔体内热辐射的谱能量密度只依赖频率与温度, 与壁材料、腔体形状无关, 且平衡时各向同性、无偏振. $u_\nu(\nu,T)\dd\nu$是单位体积在该频带内的能量; $u(T)=\int u_\nu\dd\nu$是总能量密度.

各向同性辐射落在面积元上的通量, 对入射半球积分法向速度$c\cos\theta$:
\[
j=\frac{cu}{4\pi}\int_0^{2\pi}\dd\phi\int_0^{\pi/2}\sin\theta\cos\theta\dd\theta=\frac{cu}{4}.
\]
对镜面反射, 每个光子法向动量改变$2p\cos\theta$, 再多一个通量投影因子给$P=u/3$. 吸收率$\alpha_\nu$与发射功率$e_\nu$在平衡中满足$e_\nu/\alpha_\nu=cu_\nu/4$, 即Kirchhoff定律. 理想黑体$\alpha_\nu=1$, 黑体光谱可由腔体小孔测量; 黑色外观只涉及可见频段, 不是任意频率下的完美吸收.

仅用热力学也能得到温度幂次. $U=u(T)V$和$P=u(T)/3$, 代入$U_V|_T=T P_T|_V-P$得$u=Tu'/3-u/3$, 因而$u\propto T^4$. 比例系数需要微观统计确定.

\originalsection{模式计数与Planck谱}
电磁驻波每一个独立模式相当于谐振子, 但有两个横向偏振. 周期边界条件下每个波矢占$(2\pi)^3/V$的波矢空间体积. 在半径$k$的球壳中计数得$D(k)\dd k=Vk^2\dd k/\pi^2$. 令$\omega=ck$, 则$D(\omega)=V\omega^2/(\pi^2c^3)$. 导体盒的第一卦限驻波计数给同样体相结果, 边界修正随尺寸增大相对消失.

若给每模经典均分能$\kb T$, 总能量积分$\int_0^\infty\omega^2\kb T\dd\omega$发散, 是紫外灾难. 量子热激发能为$\hbar\omega/(\e^{\beta\hbar\omega}-1)$, 因而
\[
u_\omega(\omega,T)=\frac{\hbar\omega^3}{\pi^2c^3(\e^{\hbar\omega/\kb T}-1)},\quad
u_\nu(\nu,T)=\frac{8\pi h\nu^3}{c^3(\e^{h\nu/\kb T}-1)}.
\]
这里左端记号均为$u$而非速度: $u_\nu\dd\nu=u_\omega\dd\omega$, $\omega=2\pi\nu$. 零点项$\hbar\omega/2$是温度无关的真空能, 不属于可测的热黑体谱; 完整真空能及其边界差需另作量子场论处理, 不能用本式声称已解决发散.

低频$\hbar\omega\ll\kb T$恢复Rayleigh--Jeans式, 高频由指数压低. 对$u_\omega$求导给非零谱峰方程$3(1-\e^{-x})=x$, $x=\hbar\omega/\kb T$, 根约2.82144. 若按单位波长画谱, Jacobian会把幂次改成5, 峰方程也变为$5(1-\e^{-x})=x$. 光子数谱比能量谱少一项$h\nu$, 又有不同峰位.

\originalsection{辐射热力学}
频率积分中的常数可写成Riemann（黎曼）zeta函数. 本书所需的实参数定义与收敛范围是
\[\zeta(s)=\sum_{n=1}^{\infty}\frac1{n^s},\qquad s>1.\]
与$\int_1^\infty x^{-s}\dd x$比较可知$s>1$时级数收敛, $s\le1$时发散. 对$s>1$, 展开$1/(\e^x-1)=\sum_{n\ge1}\e^{-nx}$, 各项非负, 可用单调收敛交换积分与和; 再令$y=nx$, 得
\[\int_0^\infty\frac{x^{s-1}}{\e^x-1}\dd x=\sum_{n\ge1}\frac{\Gamma(s)}{n^s}=\Gamma(s)\zeta(s).\]
原点附近被积函数约为$x^{s-2}$, 也直接说明同一收敛条件. 取$s=4$, $\Gamma(4)=3!=6$, 从而积分为$6\zeta(4)=\pi^4/15$. 后文Bose凝聚容量用到的$\zeta(3/2)$与$\zeta(3)$也都在此范围内. 得
\[
u(T)=aT^4,\quad a=\frac{\pi^2\kb^4}{15\hbar^3c^3},\qquad
j=\sigma T^4,\quad\sigma=\frac{\pi^2\kb^4}{60\hbar^3c^2}.
\]
$\sigma\simeq5.6704\times10^{-8}\,\mathrm{W\,m^{-2}K^{-4}}$. 热配分函数满足$\ln Z=-\sum_\alpha\ln(1-\e^{-\beta\hbar\omega_\alpha})$. 对模式积分, 或将Planck式积分分部得到
\[
F=-\frac13aVT^4,\qquad S=\frac43aVT^3,\qquad P=\frac13aT^4,\qquad C_V=4aVT^3.
\]
可逆绝热路径$VT^3=\text{常数}$, 也给$PV^{4/3}=\text{常数}$. 光子可以被壁吸收和产生, 总光子数不是守恒约束, 所以化学势为零. 平衡光子数$N_\gamma/V=2\zeta(3)(\kb T)^3/(\pi^2\hbar^3c^3)$, 平均光子能量$\pi^4\kb T/[30\zeta(3)]\simeq2.701\kb T$. 本稿自算的Planck曲线是理论谱, 不冒称替代原课的COBE测量数据.
